% !TEX program = xelatex
% =============================================================================
% cn-everyday · 收条（律所级要件）
%
% 收款凭证。按律师起草标准写全：付款人身份、收款方式、金额大写＋小写、
% 款项性质与对应债权（指明是哪份借条项下）、冲抵说明、尚欠余额、
% 剩余债权的权利保留声明（防"收条=结清"的常见争议）、份数。
% 样张与借条（iou.tex）同一虚构案件：2024-12-10 收到部分还款 20,000 元，
% 尚欠 180,000 元——与 cn-litigation/complaint.tex 起诉状所述一致。
%
% 编译：bash scripts/compile.sh templates/cn-everyday/receipt.tex --preview
% =============================================================================
\documentclass[zihao=-4]{ctexart}
\usepackage{everyday}

\begin{document}

\edTitle{收\hspace{1\ccwd}条}

今收到付款人\hl{王××}（身份证号：\edblank[10em]{110×××××××××××××××}）
以银行转账方式归还的借款人民币\hl{贰万元整}（小写：¥\hl{20,000.00}）。

\hl{一、款项性质。}上述款项系王××就其2024年3月1日向本人出具的
《借条》项下借款人民币贰拾万元所作的部分还款，全额冲抵借款本金。

\hl{二、余额确认。}截至本收条出具之日，该《借条》项下尚欠借款本金
人民币\hl{壹拾捌万元整}（小写：¥\hl{180,000.00}）。

\hl{三、权利保留。}本收条仅证明本人收到上述款项，不构成对剩余债权的
放弃或减免；本人对剩余本金及《借条》项下逾期利息、实现债权费用等全部
权利予以保留。

\hl{四、份数。}本收条一式两份，收款人、付款人各执一份，两份具有同等
效力。

\vspace{0.8\baselineskip}
\edSign{收款人（签名）：}{林××}
\edSign[10em]{身份证号：}{110×××××××××××××××}
\edSign[9em]{出具日期：}{2024年12月10日}

\end{document}
